% Generated draft: requires mathematical and source review.
\documentclass[11pt]{article}
\usepackage[margin=25mm]{geometry}
\usepackage{fontspec}
\setmainfont{Latin Modern Roman}
\usepackage{amsmath,amssymb,mathtools,unicode-math}
\setmathfont{Latin Modern Math}
\usepackage{enumitem,needspace,xurl,hyperref}
\setlength{\parindent}{0pt}
\setlength{\parskip}{5pt}
\setlength{\emergencystretch}{3em}
\newcommand{\N}{\mathbb{N}}
\newcommand{\Z}{\mathbb{Z}}
\newcommand{\Q}{\mathbb{Q}}
\newcommand{\R}{\mathbb{R}}
\newcommand{\C}{\mathbb{C}}
\newcommand{\F}{\mathbb{F}}
\newcommand{\A}{\mathbb{A}}
\newcommand{\PP}{\mathbb{P}}
\newcommand{\E}{\mathbb{E}}
\newcommand{\eps}{\varepsilon}
\newcommand{\dd}{\,\mathrm{d}}
\begin{document}
\section*{UnsolvedMath ID: 20001932}
\textbf{Source problem code:} catalogue entry
\subsection{Definitions and mathematical statement}
Let $g$ be a smooth Riemannian metric on $T^2=\R^2/(2\pi\Z)^2$, with geodesic vector field $X_g$ on $TT^2$ and energy $E(x,v)=\tfrac12g_x(v,v)$. A first integral is a smooth function $F$ with $X_gF=0$. A fiberwise polynomial integral of exact degree three has an expansion $F=F_0+F_1+F_2+F_3$, with $F_j(x,tv)=t^jF_j(x,v)$ and $F_3\not\equiv0$. This means exact degree, not degree at most three. Homogeneous components are separately first integrals for the geodesic flow. Thus one may state the question for a nonzero homogeneous cubic
\[
 F_3(x,v)=K_x(v,v,v),\qquad K\in\Gamma(\operatorname{Sym}^3T^*T^2),
 \qquad \nabla_{(i}K_{jkl)}=0.
\]
A Killing vector field here is a vector field $V\not\equiv0$ with $\mathcal L_Vg=0$. It gives the linear first integral $I_V(x,v)=g_x(V_x,v)$ and hence the cubic $I_V^3$.

\textbf{Conjecture 10.2.1, [Bo-Ko-Fo1995].} Every Riemannian metric on $T^2$ admitting a nonzero homogeneous cubic first integral admits a nonzero Killing vector field. Equivalently, an integral of exact degree three should force a nonzero linear integral. No global analyticity assumption is added to the metric or coefficient tensor.

\textbf{Supplied reductions and convention warning.} The research aid asserts that the degree-at-most-three version is false: for $0<\eps<1/4$,
\[
 g_\eps=\exp\bigl(2\eps(\cos x+\cos y)\bigr)(\dd x^2+\dd y^2)
\]
is an analytic torus metric with no nonzero Killing field, although its energy is a degree-two integral. This does not refute the exact-degree conjecture. It also gives the pure-trace reduction: if a cubic integral has the global form $F_3=E L$, where $L$ is fiberwise linear, then $L$ is a nonzero linear integral. The summary further claims a reduction for ``otherwise polynomially reducible'' cubics, but does not specify the global factorization hypotheses; that claim is not used to enlarge the precise pure-trace statement. The aid describes the globally irreducible, everywhere-hyperbolic branch as unresolved; its PDE terminology and hypotheses are absent.

\textbf{Further questions and historical context in the same extract.} The source attributes to Darboux and Whittaker the homogeneous-component observation for integrals analytic in velocities; this is understood as an expansion at the zero section. It attributes to [Kolokoltsov1983] nonexistence, on a closed surface of genus at least two, of a polynomial integral functionally independent of energy. For $S^2$ and $T^2$, its historical table concerns the \emph{smallest degree of a nontrivial polynomial integral}: degrees one and two have effective classifications [Bo-Ma-Fo1998]; degrees three and four have series of examples on the sphere and partial negative results on the torus; for degrees at least five it says ``nothing is known''. These are dated source descriptions, not present-day conclusions.

A related version concerns a natural Lagrangian
\[
 L(x,v)=\sum_{i,j}g_{ij}(x)v^iv^j+U(x)
\]
and a first integral with cubic and linear velocity terms. The source says most cited partial negative results assume $g$ flat, citing [Bialy1987, Mironov2010, De-Ko-Tr2012], while [Bialy2010] does not require that assumption. The displayed plus sign before $U$ is retained; it is the source's Lagrangian convention.

The source proposes analogous torus conjectures for every degree $d$: odd exact degree should force a Killing field; even degree should force a quadratic integral not proportional to energy. In the even case a nontriviality hypothesis is essential: powers $E^{d/2}$ are always integrals. A precise \emph{intended nontrivial version} assumes $\dd F\wedge\dd E\not\equiv0$ and asks for a quadratic integral $Q$ with $\dd Q\wedge\dd E\not\equiv0$. This explicitly repairs an omitted premise rather than asserting the literal even-degree wording.

The adjacent Section 10.3 asks for further two-dimensional metrics for which all geodesics can be written using elementary functions, beyond constant-curvature and Darboux-superintegrable examples [Br-Ma-Ma2008], and for metrics with explicitly given distance functions [Tao2010]. The admissible meaning of ``explicit'' is not specified. The final Section 11, communicated by Rotman, defines a stationary net in a Riemannian surface as a graph with geodesic edges, valence at least three, and $\sum_e u_e=0$ at each vertex, where $u_e$ are outgoing unit tangent vectors. It defines a $\Theta$-graph as two vertices joined by three edges. No complete stationary-net existence question occurs in this record.
\subsection{Short English statement}
The conjecture asks whether a genuine cubic geodesic integral on a two-torus forces a continuous isometric symmetry. Exact degree three matters: energy alone would satisfy a degree-at-most-three reading. The supplied aid treats a pure-trace special case and leaves the general conjecture unresolved; higher-degree and adjacent explicit-geodesic questions are retained with their missing conventions.
\subsection{Sources}
\begin{itemize}
\item Local UnsolvedMath registry, permanent ID 20001932, source code catalogue entry; canonical \texttt{aim-geometry-notes.json}, notes[269].
\item AIM Geodesics problem list, Conjecture 10.2.1 [Bo-Ko-Fo1995], with Sections 10.3 and 11 in the same extracted text: \url{https://aimath.org/pastworkshops/geodesicsproblems.pdf}.
\item Supplied labels [Kolokoltsov1983], [Bo-Ma-Fo1998], [Bialy1987], [Mironov2010], [De-Ko-Tr2012], [Bialy2010], [Br-Ma-Ma2008], [Tao2010]. Full bibliographic entries are not supplied. Historical observation attributed to Darboux and Whittaker; stationary nets communicated by Rotman.
\item Registry research attempt 1, dated 2026-08-14: partial result, status at run open, rigor ``proved'', novelty confidence low. This is an explicitly machine-generated, non-peer-reviewed aid requiring independent verification before citation.
\end{itemize}
\end{document}
